\documentclass[11pt,letterpaper]{article}
\usepackage[T1]{fontenc}
\usepackage[utf8]{inputenc}
\usepackage{amsmath,amsthm,mathtools}
\usepackage{newpxtext,newpxmath}
\usepackage[margin=1in,headheight=15pt]{geometry}
\usepackage{microtype}
\usepackage{booktabs,array,longtable}
\usepackage{xcolor}
\usepackage[most]{tcolorbox}
\usepackage{enumitem}
\usepackage{listings}
\usepackage{fancyhdr}
\usepackage{titlesec}
\usepackage{aliascnt}
\usepackage{hyperref}
\usepackage[nameinlink,noabbrev]{cleveref}
\definecolor{OliveInk}{HTML}{4C5942}
\definecolor{SoftPaper}{HTML}{F3F3E9}
\definecolor{RuleGray}{HTML}{C8CCBD}
\definecolor{MutedInk}{HTML}{52584D}
\hypersetup{colorlinks=true,linkcolor=OliveInk,citecolor=OliveInk,urlcolor=OliveInk,
 pdftitle={Ordinal complexity of graph minors},
 pdfauthor={Research manuscript prepared with ChatGPT},
 pdfsubject={Maximal order types and ordinal widths of degree-two graph-minor orders}}
\titleformat{\section}{\Large\bfseries\color{OliveInk}}{\thesection}{0.65em}{}
\titleformat{\subsection}{\large\bfseries\color{OliveInk}}{\thesubsection}{0.65em}{}
\pagestyle{fancy}
\fancyhf{}
\fancyhead[L]{\small\color{MutedInk}Ordinal complexity of graph minors}
\fancyhead[R]{\small\color{MutedInk}Research manuscript}
\fancyfoot[C]{\small\thepage}
\renewcommand{\headrulewidth}{0.3pt}
\setlength{\parskip}{0.35em}
\setlength{\parindent}{1.15em}
\setlist{itemsep=0.25em,topsep=0.35em}
\emergencystretch=2em
\theoremstyle{plain}
\newtheorem{theorem}{Theorem}[section]
\newaliascnt{lemma}{theorem}
\newtheorem{lemma}[lemma]{Lemma}
\aliascntresetthe{lemma}
\newaliascnt{proposition}{theorem}
\newtheorem{proposition}[proposition]{Proposition}
\aliascntresetthe{proposition}
\newaliascnt{corollary}{theorem}
\newtheorem{corollary}[corollary]{Corollary}
\aliascntresetthe{corollary}
\theoremstyle{definition}
\newaliascnt{definition}{theorem}
\newtheorem{definition}[definition]{Definition}
\aliascntresetthe{definition}
\newaliascnt{example}{theorem}
\newtheorem{example}[example]{Example}
\aliascntresetthe{example}
\newaliascnt{problem}{theorem}
\newtheorem{problem}[problem]{Problem}
\aliascntresetthe{problem}
\theoremstyle{remark}
\newaliascnt{remark}{theorem}
\newtheorem{remark}[remark]{Remark}
\aliascntresetthe{remark}
\newcommand{\N}{\mathbb N}
\newcommand{\Ord}{\operatorname{Ord}}
\newcommand{\otp}{\operatorname{otp}}
\newcommand{\mot}{\mathsf o}
\newcommand{\wid}{\mathsf w}
\newcommand{\hei}{\mathsf h}
\newcommand{\minor}{\preccurlyeq_{\!m}}
\newcommand{\Memb}{\mathcal M_{\mathrm{emb}}}
\newcommand{\Cl}{\mathcal C}
\newcommand{\Add}{\mathcal A}
\newcommand{\LF}{\mathcal L}
\newcommand{\DG}{\mathcal D}
\newcommand{\Q}{\mathcal Q}
\newcommand{\bigop}{\mathop{\bigoplus}}
\newcommand{\sqgraph}{\mathbin{\dot\cup}}
\newcommand{\down}{\mathord\downarrow}
\newcommand{\Inf}{\infty}
\newcommand{\eps}{\varepsilon}
\newtcolorbox{resultbox}[1][]{colback=SoftPaper,colframe=RuleGray,
 boxrule=0.5pt,arc=1mm,left=2.5mm,right=2.5mm,top=2mm,bottom=2mm,#1}
\lstset{basicstyle=\ttfamily\small,columns=fullflexible,breaklines=true,
 frame=single,rulecolor=\color{RuleGray},backgroundcolor=\color{SoftPaper},
 showstringspaces=false,aboveskip=0.8em,belowskip=0.8em}

\begin{document}
\begin{titlepage}
\vspace*{0.65in}
\noindent{\small\sffamily\color{OliveInk}ORDER THEORY \quad / \quad GRAPH MINORS \quad / \quad ORDINAL INVARIANTS}
\vspace{0.45in}
\begin{flushleft}
{\huge\bfseries Ordinal complexity\par of graph minors\par}
\vspace{0.2in}
{\Large Component splitting, degree-two graphs,\par and discontinuous limits\par}
\vspace{0.45in}
{\large A research manuscript prepared for Vladimir Reshetnikov}\par
\vspace{0.12in}
20 September 2026
\end{flushleft}
\vspace{0.4in}
\begin{abstract}
The ordinal invariants of the graph-minor ordering provide a concrete entry
point into an open program in well-partial-order theory. We study its
restriction to finite simple graphs of maximum degree two, and to related
classes with bounds on path size and the number of cyclic components. A
component-splitting argument, combined with the established maximal-order-type
theorem for finite multisets, proves that the finite disjoint-union closure of
a connected graph class of maximal order type $\alpha<\eps_0$ has maximal order
type $\omega^\alpha$. In particular, the degree-two graph-minor order has
maximal order type and ordinal width both equal to
$\omega^{\omega\cdot2}$, while its height is only $\omega$.
We give exact formulas throughout a two-parameter family, including
$\omega^{\omega+k}$ when at most $k$ cycles are allowed and path sizes are
unbounded. These formulas yield explicit failures of continuity under
increasing unions, even unions of minor-closed classes. An exact
cycle-allocation and path-packing criterion is proved and checked against
closure under elementary minor operations on all $1{,}500$ isomorphism classes
with at most fourteen vertices. The work is a partial attack on the general
graph-minor problem, not a calculation for all finite graphs; historical
novelty of the special-case formulas is not asserted.
\end{abstract}
\vfill
\begin{resultbox}
\textbf{Central calculation.}\quad For $\DG$, the finite simple graphs of
maximum degree at most two, ordered by minors,
\[
\boxed{\mot(\DG)=\wid(\DG)=\omega^{\omega\cdot2},\qquad \hei(\DG)=\omega.}
\]
The archive includes this manuscript, its \LaTeX{} source, an independent
finite checker, machine-readable results, and a proof audit.
\end{resultbox}
\end{titlepage}

\begingroup
\small
\setlength{\parskip}{0pt}
\tableofcontents
\endgroup
\clearpage

\section{The selected problem and the scope of the attack}

A well-partial-order is not determined, even approximately, by its
cardinality. The coordinatewise orders on $\N$, $\N^2$, and $\N^3$ are all
countable, have no infinite antichain, and have only finite descending
sequences. Nevertheless, their maximal linear-extension types are
$\omega$, $\omega^2$, and $\omega^3$. The ordinal invariant records how much
freedom remains when one resolves incomparabilities into a total order.

Vialard's 2024 thesis identifies the ordinal measurement of finite
undirected graphs under the minor relation as a promising unresolved
direction \cite[p.~110]{VialardThesis}. The broad question motivating this
manuscript is the following.

\begin{problem}[Graph-minor ordinal measurement]
Determine the maximal order type, and associated ordinal invariants, of
finite graphs under the minor relation, preferably in a form that interacts
with labels and structural decompositions.
\end{problem}

The present attack focuses on a class where the graph-theoretic geometry
can be described exactly, while the resulting order types are already
substantially transfinite. Every connected finite simple graph of maximum
degree at most two is either a path or a cycle. Disconnected graphs in this
class are therefore finite multisets of these components. The main obstacle
is that \emph{minor containment is not multiset embedding}: one source
component can split into several target components.

For instance,
\[
 P_2\sqgraph P_2\minor P_4,
 \qquad
 P_2\sqgraph P_3\minor C_5,
\]
where $P_n$ has $n$ vertices and $C_n$ is a simple $n$-cycle. Any argument
that requires distinct source components for distinct target paths would
miss these minor relations. Conversely,
\[
 C_3\sqgraph P_1\not\minor C_4:
\]
retaining a cycle uses the entire cyclic source component, leaving no
vertices for a separate path. Understanding both facts is essential.

\subsection{Results obtained}

For $b\in\{1,2,3,\ldots,\Inf\}$ and
$k\in\{0,1,2,\ldots,\Inf\}$, let $\DG_{b,k}$ consist of the graphs whose
path components have at most $b$ vertices and which have at most $k$
cyclic components. The symbol $\Inf$ means that the corresponding bound is
absent; it is not an ordinal parameter. Cycle lengths are unrestricted
throughout. Write $\DG=\DG_{\Inf,\Inf}$.

\begin{theorem}[Complete two-parameter calculation]\label{thm:table}
With $b\ge1$ and the preceding conventions, the invariants are as follows.
\begin{center}
\normalfont\small
\renewcommand{\arraystretch}{1.35}
\begin{tabular}{@{}lllcc@{}}
\toprule
Path-size bound & Cycle-count bound & $\mot(\DG_{b,k})$ & $\wid(\DG_{b,k})$ & $\hei$\\
\midrule
$b<\omega$ & $k<\omega$ & $\omega^{b+k}$ & $\omega^{b+k-1}$ & $\omega$\\
$b=\Inf$ & $k<\omega$ & $\omega^{\omega+k}$ & $\omega^{\omega+k}$ & $\omega$\\
$b<\omega$ & $k=\Inf$ & $\omega^\omega$ & $\omega^\omega$ & $\omega$\\
$b=\Inf$ & $k=\Inf$ & $\omega^{\omega\cdot2}$ & $\omega^{\omega\cdot2}$ & $\omega$\\
\bottomrule
\end{tabular}
\end{center}
\end{theorem}

The proof occupies \cref{sec:prelim,sec:components,sec:finite,sec:geometry,sec:unbounded,sec:transfer,sec:bounded}.
The result is not based on extrapolating finite computations.

A general component theorem is also proved: if $\Cl$ is a well-partial-order
of connected finite simple graphs with $\mot(\Cl)=\alpha<\eps_0$, then its
finite disjoint-union closure has maximal order type $\omega^\alpha$.
For a finite collection of $N\ge1$ connected templates, the corresponding
maximal order type and width are $\omega^N$ and $\omega^{N-1}$,
independently of the detailed minor relations between the templates.

\subsection{What is and is not claimed}

The standard facts about maximal linearizations, natural products,
multisets, and the height--width inequality are attributed below. The
component reduction and the degree-two calculations are given with proofs.
They are applications and extensions of that established machinery, not a
claim to have invented it. The literature search located the broad open
program, but did not establish priority for these particular special cases.
The formulas should therefore be read as explicit, auditable mathematical
results derived in this investigation, not as a claim that a named
longstanding conjecture has been resolved.

The manuscript does not determine $\mot$ for all finite graphs, for all
finite trees, or for edge-labeled graph-minor orders. It does identify where
the present method is exact, gives a nontrivial parameter spectrum, and
isolates a genuine obstruction to extending the same proof indiscriminately.

\section{Ordinal conventions and the established tools}\label{sec:prelim}

\subsection{Three different ranks}

We use $\N=\{0,1,2,\ldots\}$. All graphs are considered up to isomorphism. For finite simple graphs, the
minor relation on isomorphism classes is antisymmetric. A proper minor
operation strictly decreases $|V|+|E|$, and mutual minor containment
therefore forces isomorphism.

A sequence $(x_0,\ldots,x_{r-1})$ in a partial order $P$ is \emph{bad} if
$x_i\not\le x_j$ whenever $i<j$. Let $\operatorname{Bad}(P)$ be the tree of
finite bad sequences, including the empty sequence. Define
$\operatorname{Ant}(P)$ similarly using pairwise incomparable entries, and
$\operatorname{Dec}(P)$ using strictly descending sequences. For a
well-founded tree, the rank of a node is the supremum of the ranks of its
immediate extensions plus one. Leaves have rank zero. A partial order is a \emph{well-partial-order}
when its bad-sequence tree is well-founded; equivalently, it has neither
an infinite descending chain nor an infinite antichain.

\begin{definition}
The \emph{maximal order type} $\mot(P)$, \emph{ordinal width} $\wid(P)$,
and \emph{height} $\hei(P)$ are the ranks of the empty sequence in
$\operatorname{Bad}(P)$, $\operatorname{Ant}(P)$, and
$\operatorname{Dec}(P)$, respectively.
\end{definition}

In particular, a nonempty chain has width $1$, not $0$, and a finite
$n$-element antichain has width $n$. An ordinal width such as
$\omega^{\omega\cdot2}$ is a rank of the tree of finite antichains; it is
\emph{not} the cardinality or order type of an actual antichain.

The residual identities are
\begin{align}
\mot(P)&=\sup_{x\in P}\bigl(\mot(\{y:x\not\le y\})+1\bigr),\label{eq:motres}\\
\wid(P)&=\sup_{x\in P}\bigl(\wid(\{y:y\perp x\})+1\bigr),\label{eq:widres}\\
\hei(P)&=\sup_{x\in P}\bigl(\hei(\{y:y<x\})+1\bigr).\label{eq:heightres}
\end{align}
They follow directly by taking the first entry of a sequence.

\subsection{Natural arithmetic}

We write $\oplus$ and $\otimes$ for Hessenberg natural sum and product.
Natural sum adds the coefficients in a common Cantor normal form. Natural
product multiplies these normal forms as polynomials, using natural sum
for exponents. Thus
\[
 (\omega^\xi)\otimes(\omega^\eta)=\omega^{\xi\oplus\eta}.
\]
Ordinary ordinal addition and multiplication are written $+$ and
$\cdot$. The distinctions needed below include
\[
 b+\omega=\omega,\quad b\oplus\omega=\omega+b,
 \quad\text{and}\quad
 \omega^\omega\cdot\omega^k=\omega^{\omega+k}
 \quad(b,k<\omega).
\]
A finite natural sum of monomials with exponents strictly below $\xi$
is below $\omega^\xi$. Natural sum is strictly increasing in each
argument. An infinite additively indecomposable ordinal is an ordinal
$\omega^\xi$ with $\xi>0$; adding finitely many elements to an ordinal
strictly below it still leaves an ordinal strictly below it.

\subsection{Explicit external dependencies}

We use four established results. They are recorded here to make the proof
dependencies transparent.

\begin{theorem}[de Jongh--Parikh]\label{thm:standardmot}
A well-partial-order has a linear extension whose order type is exactly
$\mot(P)$. Moreover,
\[
 \mot(P\times Q)=\mot(P)\otimes\mot(Q),
 \qquad
 \mot(P\sqcup Q)=\mot(P)\oplus\mot(Q),
\]
where $\sqcup$ denotes a disjoint union with no cross-comparabilities.
\end{theorem}

These are the linearization and composition results of
\cite{deJonghParikh}. In particular, $\mot(\N^d)=\omega^d$ for a positive
finite $d$.

\begin{theorem}[Finite multiset maximal order type]\label{thm:multiset}
Let $P$ be a well-partial-order with $\mot(P)=\alpha<\eps_0$.
Order finite multisets by the existence of an injective matching that
maps each entry to a greater or equal entry. Then
\[
 \mot(\Memb(P))=\omega^\alpha.
\]
\end{theorem}

Here $\eps_0$ is the least nonzero ordinal satisfying
$\omega^{\eps_0}=\eps_0$. We use only this below-$\eps_0$ specialization of the multiset theorem
\cite{Weiermann,VanDerMeeren}; the general theorem has an exponent
correction. A modern presentation distinguishes carefully between
multiset embedding and multiset replacement \cite{VialardMultisets}.

\begin{theorem}[K\v r\'i\v z--Thomas]\label{thm:KT}
For every well-partial-order $P$,
\[
 \wid(P)\le\mot(P)\le\hei(P)\otimes\wid(P).
\]
\end{theorem}

This is the height--width inequality of \cite{KrizThomas}. The last
standard fact is closure of well-quasi-ordering under finite multisets;
it follows from Higman's theorem by forgetting the order of a finite
word \cite{Higman}.

For clarity, the following elementary comparison rules are proved rather
than left implicit. Passing to an induced suborder cannot increase any
of the three tree ranks. Adding comparabilities on the same underlying
set cannot increase $\mot$ or $\wid$: it only removes bad sequences and
antichains. Finally, if $P$ is partitioned into finitely many pieces
$P_i$, then
\begin{equation}\label{eq:finitecover}
 \wid(P)\le\bigop_i\wid(P_i),
 \qquad
 \mot(P)\le\bigop_i\mot(P_i).
\end{equation}
For maximal order type, remove all cross-comparabilities and apply the
disjoint-sum formula. For width, at an antichain-tree node take the natural
sum of the antichain ranks remaining in the individual pieces. Extending
the antichain lowers the rank in its chosen piece and does not increase
any other summand. This is a strict ordinal ranking with the displayed
root value, proving the upper bound. A finite cover can be made a partition
by assigning each point to its first containing piece; each resulting
piece is an induced suborder of the original covering piece.

\section{A component-splitting transfer theorem}\label{sec:components}

Let $\Cl$ be a collection of pairwise nonisomorphic, nonempty, connected
finite simple graphs. Give it the inherited minor relation. Define
\[
 \Add(\Cl)=\{C_1\sqgraph\cdots\sqgraph C_r:
          r<\omega,\ C_i\in\Cl\},
\]
including the empty graph when $r=0$. Neither $\Cl$ nor $\Add(\Cl)$ is
required to be minor-closed. The relation always means ordinary graph
minor containment in the ambient class of all finite graphs; intermediate
minor operations need not stay inside the displayed subclass.

\begin{lemma}[Allocation to source components]\label{lem:componentallocation}
Suppose $H\minor G$. Each connected component of $H$ can be assigned to a
unique connected component of $G$ by a minor model. The target components
assigned to a source component $C$ are each connected minors of $C$.
If one of them is isomorphic to $C$, it is the only target component
assigned to that source component.
\end{lemma}

\begin{proof}
A minor model assigns to each target vertex a nonempty connected branch
set in the source, with distinct branch sets disjoint. Branch sets
representing a connected target component, together with the edges
realizing its adjacency, all lie in one source component. This gives the
allocation and also a connected-minor model for each target component.

If a target component is isomorphic to $C$, it has $|V(C)|$ vertices.
Its $|V(C)|$ nonempty disjoint branch sets already use every vertex of
$C$. There is no vertex left to represent another target component.
\end{proof}

\begin{lemma}[Cantor-normal-form linearization]\label{lem:CNF}
Let $\ell:\Cl\to\alpha$ be a bijection that realizes a well-order
linear extension of the connected minor order. For a graph $G$ in
$\Add(\Cl)$, put
\begin{equation}\label{eq:generalcode}
 \Phi_\ell(G)=\bigop_{C\text{ a component of }G}\omega^{\ell(C)}.
\end{equation}
Then $\Phi_\ell$ is a bijection from $\Add(\Cl)$ onto the ordinals below
$\omega^\alpha$. Moreover,
\[
 H\minor G,\ H\not\cong G
 \quad\Longrightarrow\quad
 \Phi_\ell(H)<\Phi_\ell(G).
\]
Consequently, comparing these codes gives a linear extension of type
$\omega^\alpha$.
\end{lemma}

\begin{proof}
Unique component decomposition and uniqueness of Cantor normal form show
that the map is injective. Conversely, every ordinal below
$\omega^\alpha$ is a finite Cantor-normal-form polynomial with exponents
below $\alpha$. For each exponent, take the graph component with that
$\ell$-value and use the required finite multiplicity. This proves
surjectivity, including the empty graph and the ordinal zero.

Now allocate target components to source components as in
\cref{lem:componentallocation}. If a source component $C$ survives
unchanged, its contribution is $\omega^{\ell(C)}$. Otherwise, it either
produces no target component, or produces a finite list of proper
connected minors $D$ with $\ell(D)<\ell(C)$. Their total natural-sum
contribution is strictly smaller than $\omega^{\ell(C)}$. In the case
$\ell(C)=0$, only the empty list can occur in this latter alternative.
Thus each source contribution weakly decreases, and every changed
source contribution strictly decreases. If none changes, the whole
output is isomorphic to the input. Therefore a proper minor changes at
least one contribution. Strict monotonicity of finite natural sums
proves the asserted inequality.
\end{proof}

\begin{theorem}[Component transfer below $\eps_0$]\label{thm:component}
If $\Cl$ is a well-partial-order under minors and
$\mot(\Cl)=\alpha<\eps_0$, then $\Add(\Cl)$ is a well-partial-order and
\[
 \boxed{\mot(\Add(\Cl))=\omega^\alpha.}
\]
\end{theorem}

\begin{proof}
Identify a graph with its finite multiset of connected components.
An injective componentwise minor matching is sufficient for a graph
minor model, since the corresponding models can be placed in distinct
source components. Hence the minor order on $\Add(\Cl)$ is an
augmentation of $\Memb(\Cl)$. This proves well-partial-ordering and,
by \cref{thm:multiset}, gives the upper bound $\omega^\alpha$.

By \cref{thm:standardmot}, the connected order has a linear extension
of type $\alpha$. Apply \cref{lem:CNF} to obtain a linear extension of
$\Add(\Cl)$ of type $\omega^\alpha$. This is the matching lower bound.
\end{proof}

\begin{remark}[Why splitting does not spoil the calculation]
Componentwise matching is used only for the upper bound. The lower bound
permits a source component to split, but charges every resulting proper
connected minor below the source exponent. A finite number of lower
monomials cannot catch a higher monomial. This is the precise reason why
the extra minor comparisons do not lower the maximal order type in the
range covered by the theorem.
\end{remark}

\begin{remark}[The boundary of the theorem]\label{rem:epsilon}
The code construction is valid for any ordinal $\alpha$ for which a
connected linearization is available. Its lower bound is always
$\omega^\alpha$. The upper bound supplied by multiset embedding agrees
with it below $\eps_0$, but not uniformly at larger ordinals. Thus the
restriction on $\alpha$ is part of the argument, not a dispensable
technicality. Computing $\mot(\Cl)$ itself can also be the hard part.
\end{remark}

\section{Finite connected templates and linear forests}\label{sec:finite}

\subsection{Width of a finite-dimensional grid}

\begin{lemma}\label{lem:grid}
For every integer $d\ge1$,
\[
 \mot(\N^d)=\omega^d,\qquad
 \hei(\N^d)=\omega,\qquad
 \wid(\N^d)=\omega^{d-1}.
\]
\end{lemma}

\begin{proof}
The maximal order type follows from the natural-product theorem.
The sum of the coordinates is a finite strict rank, and there is an
infinite chain, so the height is $\omega$.

For the width upper bound, use induction on $d$. The case $d=1$ is a
chain. For $d\ge2$ and $x\in\N^d$, every point incomparable with $x$
has some coordinate $y_i<x_i$. The incomparable residual is therefore
covered by the finitely many slices
\[
 S_{i,t}=\{y:y_i=t\},\qquad 1\le i\le d,\quad 0\le t<x_i.
\]
Each slice is isomorphic to $\N^{d-1}$. By induction and
\cref{eq:finitecover}, its residual has width at most
$\omega^{d-2}\cdot(x_1+\cdots+x_d)$. Taking the supremum in
\cref{eq:widres} yields $\wid(\N^d)\le\omega^{d-1}$.

For $d\ge2$, suppose the inequality were strict. An ordinal
$\beta<\omega^{d-1}$ has only finite Cantor-normal-form exponents at
most $d-2$. The ordinal $\omega\otimes\beta$ then has leading exponent
at most $d-1$, and is strictly below $\omega^d$. This contradicts
\cref{thm:KT}. The case $d=1$ already gives equality.
\end{proof}

\subsection{A formula independent of template geometry}

\begin{theorem}[Finite template theorem]\label{thm:finite}
Let $\Cl$ contain exactly $N\ge1$ isomorphism types of nonempty connected
finite simple graphs. Then
\[
 \mot(\Add(\Cl))=\omega^N,\qquad
 \wid(\Add(\Cl))=\omega^{N-1},\qquad
 \hei(\Add(\Cl))=\omega.
\]
\end{theorem}

\begin{proof}
The finite connected order has maximal order type $N$, so
\cref{thm:component} gives $\omega^N$. There is also a direct upper
bound: record the $N$ component multiplicities. Coordinatewise increase
of this vector is sufficient for minor containment, so the graph order
is an augmentation of $\N^N$.

Every strict minor decreases $|V|+|E|$, which supplies a finite rank to
each graph. Repeating any fixed template gives an infinite chain. Thus
the height is $\omega$. The same multiplicity representation and
\cref{lem:grid} give $\wid\le\omega^{N-1}$. If $N\ge2$, a smaller width
would imply $\omega\otimes\wid<\omega^N$, contradicting
\cref{thm:KT}. If $N=1$, nonemptiness gives $\wid\ge1$, and the upper
bound is already $1$.
\end{proof}

The surprising feature is not that a finite list has a finite connected
order type, but that adding all finite disjoint unions produces precisely
one factor of $\omega$ per connected isomorphism type. Allowing
contractions, edge deletions, and component splitting changes the
comparabilities but not this answer.

\subsection{Paths alone}

Let $\LF_b$ be the minor order on finite disjoint unions of
$P_1,\ldots,P_b$, and let $\LF=\LF_{\Inf}$ allow all finite paths.

\begin{corollary}\label{cor:linearforests}
For finite $b\ge1$,
\[
 \mot(\LF_b)=\omega^b,\qquad
 \wid(\LF_b)=\omega^{b-1},\qquad
 \hei(\LF_b)=\omega.
\]
For unrestricted linear forests,
\[
 \mot(\LF)=\wid(\LF)=\omega^\omega,\qquad \hei(\LF)=\omega.
\]
\end{corollary}

\begin{proof}
The finite statements are \cref{thm:finite}. The connected path order
$P_1<P_2<\cdots$ has type $\omega$, so
\cref{thm:component} gives $\mot(\LF)=\omega^\omega$.
Each $\LF_b$ is an induced suborder of $\LF$, whence
\[
 \wid(\LF)\ge\sup_{b\ge1}\omega^{b-1}=\omega^\omega.
\]
The reverse bound is $\wid\le\mot$. The vertex--edge norm bounds
height by $\omega$, and the isolated-vertex graphs give an infinite
chain.
\end{proof}

For a linear forest with $p_n$ copies of $P_n$, a maximal linearization
is explicitly
\begin{equation}\label{eq:pathcode}
 \Phi_P(G)=\sum_{n\text{ decreasing}}\omega^{n-1}p_n.
\end{equation}
The sum is finite and written in decreasing exponent order.
It enumerates all ordinals below $\omega^\omega$.

\section{The exact minor geometry of paths and cycles}\label{sec:geometry}

\begin{lemma}[Connected comparisons]\label{lem:connecteddegree2}
For valid positive path sizes and cycle sizes at least three,
\begin{align*}
 P_a\minor P_b&\iff a\le b,\\
 C_a\minor C_b&\iff a\le b,\\
 P_a\minor C_b&\iff a\le b,\\
 C_a\minor P_b&\text{ never holds}.
\end{align*}
A source cycle that supplies a cyclic target component cannot supply
any other target component.
\end{lemma}

\begin{proof}
A minor cannot have more vertices than its source. The positive cases
are realized by shortening a path, contracting a cycle, or opening a
cycle into a path and shortening it. A path cannot yield a cycle:
deletions and contractions preserve acyclicity.

For the final assertion, retaining a cycle in a minor of $C_b$ allows
neither a vertex deletion nor an edge deletion from that component
before the cycle is formed; either deletion destroys its only cycle.
Equivalently, a cycle model in $C_b$ must cover its entire cyclic
traversal by its branch sets. All vertices of the source component
belong to the cyclic model, so there are none available for a separate
target component.
\end{proof}

\begin{lemma}[One path bin]\label{lem:pathbin}
A linear forest with component sizes $a_1,\ldots,a_r$ is a minor of
$P_n$ if and only if $\sum_i a_i\le n$. The same condition describes
the linear forests that are minors of $C_n$.
\end{lemma}

\begin{proof}
The vertex count proves necessity. For a path, choose consecutive
segments of the desired sizes, delete the edges between segments, and
delete unused vertices. Adjacent segments require an edge deletion,
not a spare separating vertex. For a cycle, first delete one edge to
obtain $P_n$, then use the path construction. The empty target is
allowed in both cases.
\end{proof}

Represent a graph by the sorted lists
\[
 G=(p_1,\ldots,p_r\ ;\ c_1,\ldots,c_s),
 \qquad p_i\ge1,\ c_j\ge3.
\]
The $p_i$ are path sizes and the $c_j$ are cycle sizes.

\begin{theorem}[Packing criterion]\label{thm:packing}
A target graph $H$ is a minor of a source graph $G$ if and only if both
of the following can be achieved together:
\begin{enumerate}[label=(\roman*)]
\item Match each target cycle to a distinct source cycle of at least
      its length. A matched source cycle is reserved in its entirety.
\item Assign each target path to a source path or to an unreserved
      source cycle, with the sum of assigned path sizes in each bin
      no greater than that source component's vertex count.
\end{enumerate}
\end{theorem}

\begin{proof}
A minor model allocates every target component to a source component.
By \cref{lem:connecteddegree2}, cycles require distinct cyclic sources
and consume them. The remaining target components are paths, and the
vertex counts in each source component give the bin constraints.
This proves necessity. Conversely, contract each reserved cycle to
its matched target cycle, and use \cref{lem:pathbin} independently in
every remaining source component. Their disjoint union realizes the
entire target graph.
\end{proof}

Total vertex capacity alone does not suffice. For example,
$P_3\sqgraph P_3$ is not a minor of $P_4\sqgraph P_2$, even though both
have six vertices. The two length-three path components cannot be
placed into the bins of sizes four and two.

\begin{proposition}[Greedy cycle reservation is exact]\label{prop:greedy}
Process the target cycles in decreasing length order, always reserving
the smallest remaining source cycle that is long enough. Then apply
\emph{exact} bin packing to the target paths using the remaining bins.
The result accepts exactly the minor pairs. If no eligible cycle exists
at some step, the minor test fails.
\end{proposition}

\begin{proof}
Consider a feasible allocation and the largest as-yet unmatched target
cycle, of length $t$. Let $f$ be the smallest available source cycle
of length at least $t$, and let $c$ be its source in the feasible
allocation. Then $f\le c$ in vertex count.

If $f$ was unused by the cyclic targets, move the target cycle from
$c$ to $f$ and move every path assigned to $f$ into $c$. This is valid
because $c$ has at least as much capacity. If $f$ was assigned to a
different target cycle of length $u$, swap the two assignments. Since
$t$ is largest, $u\le t\le c$, so the swap is valid. Thus a feasible
allocation can always be changed to use the greedy reservation.
Induction proves the claim. Nothing here says that greedy \emph{path}
packing is correct; the path stage remains an exact search.
\end{proof}

\section{The unrestricted degree-two order}\label{sec:unbounded}

Let $\Cl_2$ be the connected graphs of maximum degree at most two.
It consists of the path chain and the cycle chain, with the cross
comparisons in \cref{lem:connecteddegree2}.

\begin{lemma}\label{lem:connectedmot}
The connected minor order satisfies $\mot(\Cl_2)=\omega\cdot2$.
\end{lemma}

\begin{proof}
Removing every cross-comparison gives the disjoint sum of two chains
of type $\omega$. Its maximal order type is
$\omega\oplus\omega=\omega\cdot2$, which is an upper bound.
The order that puts all paths first in increasing size, followed by
all cycles in increasing size, is a linear extension. Its order type
is $\omega+\omega$, giving equality.
\end{proof}

\begin{theorem}\label{thm:fullmot}
For all finite simple degree-two graphs,
\[
 \mot(\DG)=\omega^{\omega\cdot2}.
\]
A maximal linearization is supplied by the explicit code
\begin{equation}\label{eq:degreecode}
 \Phi(G)=
 \sum_{n\text{ decreasing},\ n\ge3}\omega^{\omega+(n-3)}c_n
 \; +\!
 \sum_{n\text{ decreasing},\ n\ge1}\omega^{n-1}p_n,
\end{equation}
where $c_n$ and $p_n$ denote component multiplicities.
\end{theorem}

\begin{proof}
Apply \cref{thm:component} to \cref{lem:connectedmot}.
The connected linearization used there labels $P_n$ by $n-1$ and
$C_n$ by $\omega+(n-3)$. Equation~\eqref{eq:degreecode} is precisely
\cref{eq:generalcode}. Its possible exponents are all ordinals below
$\omega\cdot2$, and all finite coefficients are allowed. It therefore
ranges bijectively over the ordinals below $\omega^{\omega\cdot2}$.
\end{proof}

For example,
\[
 G=C_5\sqgraph C_3\sqgraph C_3\sqgraph P_4\sqgraph P_1
 \quad\Longrightarrow\quad
 \Phi(G)=\omega^{\omega+2}+\omega^\omega\cdot2+\omega^3+1.
\]
Deleting the cycle edge that turns $C_5$ into $P_5$ replaces
$\omega^{\omega+2}$ by $\omega^4$, possibly adding it to existing
path terms. That is a strict ordinal decrease, regardless of the
finite multiplicities of the lower terms.

\begin{proposition}[Height and width]\label{prop:fullinvariants}
The full degree-two order has
\[
 \hei(\DG)=\omega,\qquad \wid(\DG)=\omega^{\omega\cdot2}.
\]
In fact, the ordinary well-founded rank of an individual $G\in\DG$ is
exactly $|V(G)|+|E(G)|$.
\end{proposition}

\begin{proof}
Every elementary proper minor operation strictly lowers $|V|+|E|$.
Conversely, delete the edges one at a time and then delete the isolated
vertices one at a time. All intermediate graphs still have degree at
most two. This realizes a descending chain with exactly $|V|+|E|$
steps. Thus the point rank is the asserted integer, and the height of
the entire order is $\omega$.

Put $\lambda=\omega\cdot2$. If $\beta<\omega^\lambda$, each exponent
$\xi$ in the Cantor normal form of $\beta$ is either finite or has
the form $\omega+n$ for finite $n$. In either case
$1\oplus\xi<\lambda$. Consequently
\[
 \omega\otimes\beta<\omega^\lambda.
\]
If $\wid(\DG)<\omega^\lambda$, the height--width inequality would
therefore contradict \cref{thm:fullmot}. Hence the width equals the
maximal order type.
\end{proof}

Notice that the very large code $\Phi(G)$ is \emph{not} the point rank
of $G$. It describes the position of $G$ in a particular maximal
linear extension. The intrinsic point rank is the finite integer
$|V|+|E|$. Confusing these two rankings would erase the distinction
between height and maximal order type.

\section{A width-amplification lemma}\label{sec:transfer}

We need a little more than the height--width inequality to obtain
sharp widths when only finitely many cycles are permitted. A
transferability argument supplies the missing factor of $\omega$.
The argument below is a direct tree-rank proof of the standard
transferable-order lower bound; compare \cite[Section~4.4.2]{DSS}.

For a finite subset $F$ of $P$, write
\[
 P_F=P\setminus\bigcup_{x\in F}\down x.
\]

\begin{definition}
A well-partial-order $P$ is \emph{transferable for width} if
$\wid(P_F)=\wid(P)$ for every finite $F\subseteq P$.
\end{definition}

\begin{lemma}[Finite deletion]\label{lem:finitedeletion}
Suppose every principal lower set of $P$ is finite and
$\wid(P)=\gamma$ is infinite and additively indecomposable. Then $P$
is transferable for width.
\end{lemma}

\begin{proof}
Removing the lower sets of finitely many points removes only a finite
set $E$. By \cref{eq:finitecover},
\[
 \gamma=\wid(P)\le\wid(P\setminus E)\oplus |E|.
\]
If $\wid(P\setminus E)<\gamma$, the right side is still strictly
below $\gamma$, a contradiction. The reverse inequality follows
from induced-suborder monotonicity.
\end{proof}

\begin{lemma}[Transferable product bound]\label{lem:transferbound}
If $P$ is transferable for width and $\wid(P)=\gamma$, then for every
ordinal $\beta$,
\[
 \wid(P\times\beta)\ge\gamma\cdot\beta.
\]
Consequently, for every well-partial-order $Q$,
\[
 \wid(P\times Q)\ge\gamma\cdot\mot(Q).
\]
\end{lemma}

\begin{proof}
Prove the first assertion simultaneously for $P_F$, for every finite
$F$, by transfinite induction on $\beta$. Further finite lower-set
removals from $P_F$ have the form $P_{F\cup F'}$, so they retain width
$\gamma$. The assertion is immediate for $\beta=0$. At a limit,
restrict to the suborders $P_F\times\delta$ for $\delta<\beta$ and
use continuity of multiplication in its right argument:
$\sup_{\delta<\beta}\gamma\cdot\delta=\gamma\cdot\beta$.

For the successor step, put $R=P_F$ and consider the top slice
$R\times\{\delta\}$ of $R\times(\delta+1)$. After any finite
antichain $s$ in that slice, let $A$ be its set of first coordinates.
Every point of
\[
 (R\setminus\down A)\times\delta
\]
is incomparable with all entries of $s$. Indeed, a lower-slice point
cannot lie above a top-slice point, and removing $\down A$ prevents
it from lying below one. By the induction hypothesis, the tree of
these lower-slice continuations has rank at least $\gamma\cdot\delta$.

Let $r_R(s)$ be the rank of $s$ in the antichain tree of the top
slice alone. Well-founded induction over that tree shows that the
rank of $s$ in the full product tree is at least
\[
 \gamma\cdot\delta+r_R(s).
\]
For a leaf, use the lower-slice continuation bound. At a nonleaf,
take the supremum of the same bound at each top-slice child plus one;
continuity of ordinal addition in its right argument gives the
claim. At the empty node this is
$\gamma\cdot\delta+\gamma=\gamma\cdot(\delta+1)$.

Finally, give $Q$ a maximal linear extension of type $\mot(Q)$.
The product with this linear extension is an augmentation of
$P\times Q$, so its width is no larger. The first assertion supplies
the desired lower bound.
\end{proof}

\begin{corollary}\label{cor:LFtransfer}
For any well-partial-order $Q$ with $\mot(Q)=\omega^k$, where
$k<\omega$,
\[
 \wid(\LF\times Q)=\mot(\LF\times Q)=\omega^{\omega+k}.
\]
\end{corollary}

\begin{proof}
A finite graph has only finitely many minors up to isomorphism, so
principal lower sets in $\LF$ are finite. By
\cref{cor:linearforests,lem:finitedeletion}, $\LF$ is transferable and
has width $\omega^\omega$. Therefore
\[
 \wid(\LF\times Q)\ge\omega^\omega\cdot\omega^k
       =\omega^{\omega+k}.
\]
The maximal order type of the product is
$\omega^\omega\otimes\omega^k=\omega^{\omega\oplus k}
=\omega^{\omega+k}$, which is the matching upper bound.
\end{proof}

\section{Cycle budgets and the complete parameter spectrum}\label{sec:bounded}

\subsection{A bounded number of cycles}

Let $\Q_{\le k}$ be the minor order on disjoint unions of at most
$k$ cycles, and $\Q_{=k}$ the suborder with exactly $k$ cycles. A
cycle cannot be split into two cycles, so these are finite-multiset
embedding orders over the chain of cycle lengths.

\begin{lemma}\label{lem:cyclebudget}
For $k<\omega$,
\[
 \mot(\Q_{\le k})=\mot(\Q_{=k})=\omega^k.
\]
For $k=0$, both orders are singletons, so this means $1$.
\end{lemma}

\begin{proof}
Replace a cycle length $n$ by $n-2\ge1$. Pad with zeros and sort to
obtain a nondecreasing $k$-tuple. The multiset embeds into another
exactly when its padded coordinates are componentwise no larger.
One way to see necessity is to match largest entries first: the
$j$ largest target entries need $j$ source entries at least that
large. Sufficiency follows by matching the correspondingly sorted
coordinates.

Thus $\Q_{\le k}$ is an induced suborder of $\N^k$, and has maximal
order type at most $\omega^k$. For exactly $k$ entries, lexicographic
order by the smallest coordinate first, then the next, is a linear
extension of type $\omega^k$. This is proved inductively: fixing
the first positive coordinate leaves a translated nondecreasing
$(k-1)$-tuple, of lexicographic type $\omega^{k-1}$, and there are
$\omega$ possible first coordinates. Hence the type is
$\omega^{k-1}\cdot\omega=\omega^k$. The lower bound applies to
both orders.
\end{proof}

\begin{lemma}[The exact-cycle layer is a true product]\label{lem:fixedlayer}
Let $\DG_{b,=k}$ denote the degree-two graphs with path-size bound
$b$ and exactly $k$ cycles. For finite $k$,
\[
 \DG_{b,=k}\cong\LF_b\times\Q_{=k}.
\]
On the entire class $\DG_{b,k}$, the minor order is an augmentation
of the product order $\LF_b\times\Q_{\le k}$.
\end{lemma}

\begin{proof}
For source and target with the same number of cycles, the injection
matching target cycles to source cycles is a bijection. Every source
cycle is consumed by its cyclic target and is unavailable for path
components. Thus the path part must be a minor of the source path
part, independently of the cycle comparison. This proves both
necessity and sufficiency for the exact-cycle layer.

For at-most-$k$ graphs, independent comparison of the two parts is
still sufficient. There can be additional minor comparisons because
unmatched source cycles may be opened to supply paths. This is
exactly an augmentation, rather than an isomorphism in general.
\end{proof}

\begin{theorem}[Finite cycle budgets]\label{thm:finitebudget}
For $b\ge1$ finite and $k<\omega$,
\[
 \mot(\DG_{b,k})=\omega^{b+k},\qquad
 \wid(\DG_{b,k})=\omega^{b+k-1},\qquad
 \hei(\DG_{b,k})=\omega.
\]
For unrestricted path sizes and finite $k$,
\[
 \mot(\DG_{\Inf,k})=\wid(\DG_{\Inf,k})=\omega^{\omega+k},
 \qquad \hei(\DG_{\Inf,k})=\omega.
\]
\end{theorem}

\begin{proof}
Let $\alpha_b=b$ for finite $b$ and $\alpha_{\Inf}=\omega$.
By \cref{cor:linearforests}, $\mot(\LF_b)=\omega^{\alpha_b}$.
The product upper bound from \cref{lem:fixedlayer} and the induced
exact-cycle product lower bound give
\[
 \mot(\DG_{b,k})
 =\omega^{\alpha_b}\otimes\omega^k
 =\omega^{\alpha_b\oplus k}.
\]
These are precisely the two asserted maximal order types.

For finite $b$, encode the path part by its $b$ multiplicities and
the cycle part by its padded $k$-tuple. This realizes a weaker order
as an induced suborder of $\N^{b+k}$. Consequently
$\wid(\DG_{b,k})\le\omega^{b+k-1}$. Height is $\omega$: the
vertex--edge norm gives finite point ranks, and arbitrarily many
isolated vertices give an infinite chain. The height--width
inequality forces the matching lower bound, just as in
\cref{thm:finite}; if $b+k=1$, nonemptiness supplies the lower bound
$1$ directly.

For unrestricted path sizes, apply \cref{cor:LFtransfer} to the
induced exact-cycle layer. Its width is $\omega^{\omega+k}$, already
the maximal order type of the whole at-most-$k$ class. The height
argument is unchanged.
\end{proof}

The finite-$b$ classes need not be minor-closed: opening a long
cycle can produce a path exceeding the chosen path-size bound.
The theorem concerns their inherited minor order, not a different
relation that restricts intermediate operations. This distinction
is why the proof used induced suborders and augmentations rather
than assuming closure.

\subsection{Removing the cycle-count bound}

\begin{theorem}\label{thm:unlimitedbudget}
For finite $b\ge1$,
\[
 \mot(\DG_{b,\Inf})=\wid(\DG_{b,\Inf})=\omega^\omega,
 \qquad\hei(\DG_{b,\Inf})=\omega.
\]
For unrestricted $b$, the answer is
$\mot(\DG)=\wid(\DG)=\omega^{\omega\cdot2}$ and $\hei(\DG)=\omega$.
\end{theorem}

\begin{proof}
When $b$ is finite, the connected alphabet consists of finitely many
paths and the chain of all cycles. Every path in the alphabet is a
minor of every sufficiently long cycle. In any linear extension,
all the finitely many paths therefore occur within a finite initial
portion of the cycle chain. The total order type is $\omega$.
The component theorem gives $\mot(\DG_{b,\Inf})=\omega^\omega$.
On the other hand, the finite-budget suborders give
\[
 \wid(\DG_{b,\Inf})\ge\sup_{k<\omega}\omega^{b+k-1}
 =\omega^\omega.
\]
The height remains $\omega$.

For unrestricted paths, \cref{thm:fullmot,prop:fullinvariants} have
already proved the result. There is also an independent width lower
bound from \cref{thm:finitebudget}:
\[
 \wid(\DG)\ge\sup_{k<\omega}\omega^{\omega+k}
 =\omega^{\omega\cdot2}.
\]
Only monotonicity is used in these lower bounds; equality with an
increasing-union supremum is not presumed.
\end{proof}

This completes the proof of \cref{thm:table}.

\section{Two failures of continuity}\label{sec:limits}

The formulas reveal why a finite enumeration, however extensive,
is a poor guide to a transfinite order type unless supported by a
structural argument.

\subsection{An increasing union of minor-closed classes}

Let $\mathcal K_m$ contain the degree-two graphs all of whose
connected components have at most $m$ vertices. Unlike
$\DG_{b,k}$ for finite $b$ and unbounded cycle lengths, each
$\mathcal K_m$ really is minor-closed. Contractions and deletions
cannot enlarge a connected component.

For $m=1$, its connected alphabet contains just $P_1$.
For $m\ge2$, it contains the $m$ paths and the $m-2$ cycles of sizes
at most $m$, for a total of $2m-2$ templates. Thus
\begin{equation}\label{eq:componentbound}
 \mot(\mathcal K_m)=\omega^{2m-2},\qquad
 \wid(\mathcal K_m)=\omega^{2m-3}
 \quad(m\ge2).
\end{equation}
The finite template theorem handles $m=1$ separately.

\begin{corollary}[Noncontinuity even along minor-closed subclasses]\label{cor:noncontinuity}
The classes $\mathcal K_m$ are increasing and
$\DG=\bigcup_{m\ge1}\mathcal K_m$, but
\[
 \sup_m\mot(\mathcal K_m)=\omega^\omega
 <\omega^{\omega\cdot2}=\mot\left(\bigcup_m\mathcal K_m\right).
\]
The same strict inequality holds with $\wid$ in place of $\mot$.
\end{corollary}

\begin{proof}
Every finite graph has a finite maximum component size, so the union
is $\DG$. Take the suprema in \cref{eq:componentbound} and compare
with \cref{thm:fullmot,prop:fullinvariants}.
\end{proof}

The familiar continuity of ordinal exponentiation,
$\sup_n\omega^n=\omega^\omega$, does not imply continuity of the
operation $P\mapsto\mot(P)$ on increasing unions of orders. These
are different operations. Although every finite bad sequence occurs
in some $\mathcal K_m$, its continuation possibilities in the full
order need not be bounded by the rank it has in that one subclass.

\subsection{A discontinuity already with one cycle}

Fix a finite $k\ge1$. Then
\[
 \DG_{\Inf,k}=\bigcup_{b\ge1}\DG_{b,k},
 \qquad
 \sup_b\mot(\DG_{b,k})=\omega^\omega,
\]
whereas
\[
 \mot(\DG_{\Inf,k})=\omega^{\omega+k}>\omega^\omega.
\]
The same conclusion holds for ordinal width. One allowed cycle
already produces the first strict jump, to $\omega^{\omega+1}$.
The exact-cycle layer explains its source: an unrestricted linear
forest, of type $\omega^\omega$, varies independently alongside an
unbounded cycle length, contributing another factor of $\omega$.

\subsection{A maximal code need not remain maximal on a suborder}

The explicit code \cref{eq:degreecode} is maximal on $\DG$, but
restricting it to an induced subclass can lose maximality.

\begin{proposition}\label{prop:nonmaxcode}
For finite $k\ge1$, the order induced by the code $\Phi$ on
$\DG_{\Inf,k}$ has type $\omega^{\omega+1}$. For $k\ge2$, this is
strictly below $\mot(\DG_{\Inf,k})=\omega^{\omega+k}$.
\end{proposition}

\begin{proof}
There are only finitely many multisets of at most $k$ cycles whose
largest cycle length is bounded by any given integer. The cycle
polynomials, ordered by their Cantor normal forms, consequently
have order type $\omega$: the largest length is primary, and there
are only finitely many preceding patterns for each length bound.
The empty pattern is included and does not change that type.
For each fixed cycle pattern, the unrestricted path part has type
$\omega^\omega$, and every graph in one pattern precedes every graph
in the next. The total type is therefore
$\omega^\omega\cdot\omega=\omega^{\omega+1}$.
\end{proof}

For two cycles, sorting by the largest cycle first creates only
finite cycle-pattern blocks, whereas a maximal linearization of the
exact-two-cycle order sorts by the \emph{smallest} cycle first and
has type $\omega^2$. This concrete contrast is a useful diagnostic:
a correct strictly monotone ordinal code proves a lower bound on
maximal order type, but does not automatically prove that the bound
is optimal.

\section{Exact computation and reproducible checks}\label{sec:computation}

\subsection{What was checked}

The archive contains a dependency-free Python program,
\texttt{code/verify\_degree\_two.py}. Its principal test enumerates
\emph{every} degree-two graph up to isomorphism with at most fourteen
vertices. Isomorphism is exact here: a sorted multiset of path sizes
and a sorted multiset of cycle sizes uniquely determine the graph.
No graph-isomorphism heuristic is used.

Two substantially different constructions of the minor relation are
compared. The first takes transitive closure under elementary minor
operations on component signatures. The second implements the
cycle-reservation and bin-packing criterion. A third version replaces
exhaustive cycle reservation by the proved greedy reservation rule;
it shares the exact path-packing solver with the second version.

The completed run produced the following results.
\begin{center}
\renewcommand{\arraystretch}{1.2}
\begin{tabular}{@{}lr@{}}
\toprule
Quantity & Verified count\\
\midrule
Graphs up to isomorphism, at most 14 vertices & $1{,}500$\\
Ordered source--target pairs compared & $2{,}250{,}000$\\
Minor pairs, including equality & $378{,}324$\\
Proper minor pairs with strict code decrease & $376{,}824$\\
Equal-cycle-count product comparisons & $732{,}568$\\
Individual finite height-rank checks & $1{,}500$\\
Targeted regression examples & $9$\\
\bottomrule
\end{tabular}
\end{center}

All comparisons passed. The saved report records Python 3.13.5 and a
runtime of approximately 10.1 seconds in the execution environment.
Runtime is not intended as a portable benchmark. The raw counts and
machine-readable report are in \texttt{results/counts.csv} and
\texttt{results/verification.json}.

\subsection{The elementary-operation implementation}

For a path $P_n$, deleting an edge creates
$P_i\sqgraph P_{n-i}$ for $1\le i<n$; deleting a vertex creates
$P_i\sqgraph P_{n-1-i}$, omitting zero-size pieces; contracting an
edge creates $P_{n-1}$. Every possible position is included.

For $C_n$, deleting an edge creates $P_n$, deleting a vertex creates
$P_{n-1}$, and contracting an edge creates $C_{n-1}$ when $n\ge4$.
Contracting an edge of $C_3$ creates $P_2$, because parallel edges
are suppressed in the simple-graph convention. This last case is
explicitly tested; treating it as a two-cycle would silently change
the mathematical category.

Every elementary successor has a smaller value of $|V|+|E|$.
Memoized closure therefore terminates. Since elementary operations
generate the minor relation, this construction does not assume the
packing criterion it is checking.

\subsection{The packing implementation}

The exact path-packing routine processes target paths in decreasing
size. At each step it tries every distinct residual bin capacity
large enough for the next path, subtracts the path size, sorts the
remaining capacities, and memoizes the resulting state. Equal
capacities are symmetric and need only one representative branch.
It rejects immediately when total capacity or the largest capacity
is insufficient. These are necessary tests only; remaining cases
are solved by exhaustive branching, not by a greedy heuristic.

The following pseudocode describes the greedy-cycle variant.
\begin{lstlisting}
minor(target, source):
    available_cycles = source.cycles sorted increasingly
    for need in target.cycles sorted decreasingly:
        reserve the smallest available cycle with size >= need
        if none exists: return false
        remove that cycle from available_cycles
    bins = source.paths + available_cycles
    return exact_bin_pack(target.paths, bins)
\end{lstlisting}

The program also verifies the exact-cycle product identity directly
on every tested pair with the same number of cycles. This is a
separate structural assertion, not just another expression of the
ordinal-code monotonicity test.

\subsection{Exact ordinal comparisons, not numeric approximations}

For the code $\Phi$, a path exponent $n-1$ is stored as $(0,n-1)$,
and a cycle exponent $\omega+(n-3)$ as $(1,n-3)$. A graph code is
a descending list of pairs consisting of an exponent and a positive
coefficient. Lexicographic comparison of these finite lists is
exactly comparison of the represented Cantor normal forms.

The program checks that all graph codes are distinct and that every
proper minor pair has a strict decrease. It does not assign a large
floating-point number to $\omega$, and it does not attempt to
approximate transfinite arithmetic numerically.

\subsection{An independent enumeration check}

Let $a_n$ be the number of isomorphism classes on exactly $n$ vertices.
Component multiplicities give the formal power series
\begin{equation}\label{eq:gf}
 \sum_{n\ge0}a_nx^n
 =\prod_{j\ge1}\frac{1}{1-x^j}
  \prod_{j\ge3}\frac{1}{1-x^j}.
\end{equation}
The first product chooses paths and the second cycles. A coefficient
calculation by repeated unbounded-knapsack updates is checked against
the partition-based graph enumeration. The resulting sequence through
$n=14$ is
\[
\begin{gathered}
1,1,2,4,7,11,19,29,46,70,106,156,232,334,482.
\end{gathered}
\]
These coefficients sum to $1{,}500$. Equation~\eqref{eq:gf} is an
enumeration identity derived here from the component description;
it is not used to infer any of the transfinite ranks.

\subsection{What the finite checks do not prove}

For a fixed vertex cutoff, the graph order is finite, and its maximal
order type is simply the number of graph isomorphism classes below
that cutoff. A test on fourteen vertices cannot observe
$\omega^{\omega\cdot2}$. The program checks the graph-theoretic
criterion, its implementation, the strict ordinal-code comparison,
and the exact-cycle decomposition on a finite domain. The infinite
ordinal statements depend on the written arguments and the named
standard theorems, not on computational extrapolation.

No theorem prover was used for the infinite proofs. The manuscript
and checker are mutually supporting artifacts, not a formal
verification certificate.

\section{What the attack reveals about the general problem}\label{sec:outlook}

\subsection{Disjoint unions are not always the difficult step}

Suppose a connected graph class has a known maximal order type
below $\eps_0$. The component theorem determines the maximal order
type of all its finite disjoint unions without solving the exact
minor-allocation problem between arbitrary disconnected graphs.
The upper bound ignores splitting; the lower bound tolerates every
finite split into proper connected minors. Their agreement makes
the detailed splitting geometry irrelevant to this one invariant.

This observation also explains why a difficult finite decision
procedure can coexist with a simple transfinite formula. Determining
whether a particular collection of paths fits into particular source
components is a packing problem. Determining the maximal order type
of the entire family can instead be done by a multiset upper bound
and a Cantor-normal-form lower bound.

\subsection{A concrete next step: a bounded branching template}

A natural continuation is to replace cycles by a class of connected
trees with a small, fixed branching pattern but unbounded arm lengths.
The immediate task would be to identify the connected minor order,
including all minor operations that merge or remove the distinguished
branching structure, and compute its maximal order type. If that
connected type were some $\alpha<\eps_0$, the disjoint-union answer
would then be $\omega^\alpha$ by the proved transfer theorem.

This is a conditional route, not a claimed calculation for such tree
classes. The degree-two work identifies two features that would need
to be replaced. First, a cycle retained as a cycle consumes the entire
source component, creating exact product layers. A branching tree
may retain its core while releasing other vertices for extra
components. Second, paths and cycles have an elementary capacity
description. For trees with branching, total vertex capacity alone
cannot encode branch-set geometry.

\subsection{Why the full graph-minor question remains outside this result}

The component theorem does not calculate the connected order type
for arbitrary graphs. At large ordinal parameters, its lower and
upper bounds also need not agree, as noted in
\cref{rem:epsilon}. Labeling changes the connected alphabet and may
change which splitting operations are legal. Finally, the exact
product argument for a fixed cycle budget relies on a property that
general connected graphs do not possess.

The outcome of the present attempt is therefore an exact calculation
on a nontrivial part of the graph-minor order, a general transfer
mechanism in a stated ordinal range, and explicit counterexamples
to a tempting continuity principle. These statements are the proved
content. The general ordinal measurement program remains a separate
problem.

\clearpage
\appendix
\section{Proof audit and dependency map}\label{app:audit}

The following audit separates the inherited theorems from the steps
that carry the graph-specific calculation.

\begin{longtable}{@{}p{0.26\textwidth}p{0.65\textwidth}@{}}
\toprule
Claim & Dependency and critical check\\
\midrule
\endhead
Connected-to-disconnected transfer &
Uses the maximal-linearization theorem and the below-$\eps_0$
multiset formula. The graph-specific step is
\cref{lem:componentallocation}: an unchanged connected component
cannot reproduce itself plus additional components.\\[0.6em]
Finite template width &
Uses the grid width calculation and the height--width inequality.
The multiplicity map supplies only a weaker coordinatewise order;
no converse matching assertion is assumed.\\[0.6em]
Full degree-two type &
The connected type is exactly $\omega\cdot2$, not merely $\omega$.
Paths-first, cycles-second is a valid linear extension, and the
incomparable disjoint-chain order supplies its upper bound.\\[0.6em]
Full degree-two width &
The argument checks directly that
$\omega\otimes\beta<\omega^{\omega\cdot2}$ for every smaller
$\beta$. It does not incorrectly assume that
$\omega^{\omega\cdot2}$ is multiplicatively indecomposable.\\[0.6em]
Finite cycle budget &
The source and target have the same number of cycles on the
lower-bound layer. Hence every source cycle is reserved; no path
can use any of its unused-looking length.\\[0.6em]
Unbounded-path width &
Finite principal lower sets and additively indecomposable width
make $\LF$ transferable. The product lower bound is proved by
antichain-tree induction, with ordinary multiplication in the
specified order.\\[0.6em]
Union discontinuity &
Both the subclass invariants and the full-class invariant are
calculated independently. No interchange of a union and a
rank operation is used.\\[0.6em]
Computer results &
Exhaustive within the stated vertex cutoff. Elementary-operation
closure and packing are independent descriptions of minors;
the greedy and exhaustive cycle-allocation variants share the
same path-packing subroutine.\\
\bottomrule
\end{longtable}

\paragraph{Priority and verification status.}
The written proofs are not formally checked or externally refereed.
The broad source problem is documented, while historical novelty of
the special-case answers is not established. The exact finite tests
passed; their range and limits are reported rather than converted
into a claim about infinite proof verification.

\clearpage
\section{Rebuilding the artifacts}

The code uses the Python standard library only and is written for
Python 3.9 or later. From the extracted archive root, run
\begin{lstlisting}
python code/verify_degree_two.py --max-vertices 14 --out results
python build.py
\end{lstlisting}
The second command invokes \texttt{pdflatex} repeatedly and requires
a TeX installation with the packages listed in the document
preamble. The supplied PDF is already compiled. Alternatively,
\begin{lstlisting}
latexmk -pdf -interaction=nonstopmode -halt-on-error graph_minor_ordinals.tex
\end{lstlisting}
The JSON report records the execution environment and all principal
test counts. The CSV file records enumeration counts by vertex
number. The source-review notes record the scope of the literature
check and the exclusion of a lexicographic-product question that
was found to have a later published proof.

\phantomsection
\addcontentsline{toc}{section}{References}
\begin{thebibliography}{99}

\bibitem{deJonghParikh}
D.~H.~J. de Jongh and R.~Parikh.
\emph{Well-partial orderings and hierarchies}.
Indagationes Mathematicae, 39(3):195--207, 1977.
\href{https://doi.org/10.1016/1385-7258(77)90067-1}{doi:10.1016/1385-7258(77)90067-1}.

\bibitem{Higman}
G.~Higman.
\emph{Ordering by divisibility in abstract algebras}.
Proceedings of the London Mathematical Society, (3)~2:326--336, 1952.
\href{https://doi.org/10.1112/plms/s3-2.1.326}{doi:10.1112/plms/s3-2.1.326}.

\bibitem{KrizThomas}
I.~K\v r\'i\v z and R.~Thomas.
\emph{Ordinal types in Ramsey theory and well-partial-ordering theory}.
In J.~Ne\v set\v ril and V.~R\"odl, editors,
\emph{Mathematics of Ramsey Theory}, Algorithms and Combinatorics~5,
pages 57--95. Springer, 1990.
\href{https://doi.org/10.1007/978-3-642-72905-8_7}{doi:10.1007/978-3-642-72905-8\_7}.

\bibitem{Weiermann}
A.~Weiermann.
\emph{A computation of the maximal order type of the term ordering on
finite multisets}.
In \emph{Computability in Europe 2009}, LNCS~5635, pages 488--498.
Springer, 2009.
\href{https://doi.org/10.1007/978-3-642-03073-4_50}{doi:10.1007/978-3-642-03073-4\_50}.

\bibitem{VanDerMeeren}
J.~Van der Meeren, M.~Rathjen, and A.~Weiermann.
\emph{Well-partial-orderings and the big Veblen number}.
Archive for Mathematical Logic, 54(1--2):193--230, 2015.
\href{https://doi.org/10.1007/s00153-014-0408-5}{doi:10.1007/s00153-014-0408-5}.

\bibitem{DSS}
M.~D\v zamonja, S.~Schmitz, and Ph.~Schnoebelen.
\emph{On ordinal invariants in well quasi orders and finite antichain orders}.
In \emph{Well Quasi-Orders in Computation, Logic, Language and Reasoning},
Trends in Logic~53, pages 29--54. Springer, 2020.
\href{https://doi.org/10.1007/978-3-030-30229-0_2}{doi:10.1007/978-3-030-30229-0\_2}.
Accessible version: \href{https://arxiv.org/abs/1711.00428}{arXiv:1711.00428}.
Only the linearization, rank, and transferable-order results identified
in the text are used here; no lexicographic-product formula is used.

\bibitem{VialardMultisets}
I.~Vialard.
\emph{Ordinal measures of the set of finite multisets}.
In \emph{MFCS 2023}, LIPIcs~272, article~87, pages 87:1--87:15, 2023.
\href{https://doi.org/10.4230/LIPIcs.MFCS.2023.87}{doi:10.4230/LIPIcs.MFCS.2023.87}.
Preprint: \href{https://arxiv.org/abs/2302.09881}{arXiv:2302.09881}.

\bibitem{VialardThesis}
I.~Vialard.
\emph{Measuring well-quasi orders and complexity of verification}.
Doctoral thesis, Universit\'e Paris-Saclay, 2024.
The graph-minor research direction appears on printed page~110.
\url{https://isavialard.github.io/home/mwqo.pdf}.

\end{thebibliography}
\end{document}
